% AMC 12A 2023 Problem 16 Strategic Analysis
% Generated using Strategic Problem Solving Framework v2.1

\documentclass[11pt]{article}
\usepackage{amsmath, amssymb, amsthm}
\usepackage{geometry}
\usepackage{enumitem}
\usepackage{tcolorbox}
\tcbuselibrary{skins,breakable}
\usepackage{xcolor}

% Page setup
\geometry{margin=1in}
\setlength{\parindent}{0pt}
\setlength{\parskip}{6pt}

% Custom colored boxes
\newtcolorbox{problembox}[1][]{
    colback=blue!5!white,
    colframe=blue!75!black,
    title=Problem Configuration Analysis,
    breakable,
    #1
}

\newtcolorbox{strategybox}[1][]{
    colback=pink!5!white,
    colframe=pink!75!black,
    title=Strategic Framework Application,
    breakable,
    #1
}

\newtcolorbox{tacticalbox}[1][]{
    colback=green!5!white,
    colframe=green!75!black,
    title=Tactical Selection,
    breakable,
    #1
}

\newtcolorbox{conversionbox}[1][]{
    colback=yellow!5!white,
    colframe=yellow!75!black,
    title=Conversion Techniques,
    breakable,
    #1
}

\newtcolorbox{verificationbox}[1][]{
    colback=orange!5!white,
    colframe=orange!75!black,
    title=Verification Strategy,
    breakable,
    #1
}

\newtcolorbox{roadmapbox}[1][]{
    colback=purple!5!white,
    colframe=purple!75!black,
    title=Solution Roadmap,
    breakable,
    #1
}

\newtcolorbox{competitionbox}[1][]{
    colback=red!5!white,
    colframe=red!75!black,
    title=Competition Strategy,
    breakable,
    #1
}

\newtcolorbox{pitfallbox}[1][]{
    colback=gray!5!white,
    colframe=gray!75!black,
    title=Common Pitfalls and Warnings,
    breakable,
    #1
}

\newtcolorbox{solutionbox}[1][]{
    colback=cyan!5!white,
    colframe=cyan!75!black,
    title=Detailed Solution,
    breakable,
    #1
}

\newtcolorbox{keyinsightbox}[1][]{
    colback=magenta!5!white,
    colframe=magenta!75!black,
    title=Key Insights,
    breakable,
    #1
}

\begin{document}

\title{\textbf{AMC 12A 2023 Problem 16}\\Strategic Analysis}
\author{Using Framework v2.1}
\date{}
\maketitle

\section*{Problem Statement}

Consider the set of complex numbers $z$ satisfying $|1+z+z^2|=4$. The maximum value of the imaginary part of $z$ can be written in the form $\frac{\sqrt{m}}{n}$, where $m$ and $n$ are relatively prime positive integers. What is $m+n$?

\[\textbf{(A) } 20 \qquad \textbf{(B) } 21 \qquad \textbf{(C) } 22 \qquad \textbf{(D) } 23 \qquad \textbf{(E) } 24\]

\begin{problembox}
\section*{1. Problem Configuration Analysis}

\subsection*{A. Problem Classification}
\begin{itemize}[leftmargin=*]
    \item \textbf{Competition}: AMC 12A 2023
    \item \textbf{Problem Number}: 16
    \item \textbf{Difficulty Band}: Upper-mid range (Problem 16/25)
    \item \textbf{Expected Time}: 4-5 minutes
    \item \textbf{Point Value}: 6 points (standard AMC 12 scoring)
\end{itemize}

\subsection*{B. Mathematical Domain Classification}
\begin{itemize}[leftmargin=*]
    \item \textbf{Primary Domain}: Complex Numbers
    \item \textbf{Secondary Domain}: Optimization
    \item \textbf{Tertiary Domain}: Algebra (Quadratic Formula) / Geometry (Locus)
    \item \textbf{Cross-Domain Potential}: Very High (complex numbers, algebra, optimization, geometry)
\end{itemize}

\subsection*{C. Problem Type}
\begin{itemize}[leftmargin=*]
    \item \textbf{Format}: Multiple choice (integer answer)
    \item \textbf{Question Type}: Optimization problem (maximize imaginary part)
    \item \textbf{Core Task}: Find maximum $\text{Im}(z)$ subject to constraint
    \item \textbf{Answer Form}: Sum $m + n$ where max imaginary part is $\frac{\sqrt{m}}{n}$
\end{itemize}

\subsection*{D. Given Information}
\begin{itemize}[leftmargin=*]
    \item Constraint: $|1+z+z^2| = 4$
    \item Variable: $z$ is a complex number
    \item Objective: Maximize $\text{Im}(z)$
    \item Answer format: $\text{Im}(z)_{\text{max}} = \frac{\sqrt{m}}{n}$ with $\gcd(m,n) = 1$
    \item Need to find: $m + n$
\end{itemize}

\subsection*{E. Constraints}
\begin{itemize}[leftmargin=*]
    \item $z$ must be a complex number
    \item Must satisfy the magnitude constraint exactly
    \item $m$ and $n$ must be relatively prime positive integers
    \item Answer is an integer between 20 and 24
\end{itemize}

\subsection*{F. Objective}
\begin{itemize}[leftmargin=*]
    \item \textbf{Find}: Maximum imaginary part of $z$
    \item \textbf{Express}: In form $\frac{\sqrt{m}}{n}$
    \item \textbf{Output}: Calculate $m + n$
\end{itemize}

\subsection*{G. Key Structural Features}
\begin{itemize}[leftmargin=*]
    \item \textbf{Magnitude constraint}: $1+z+z^2$ lies on circle of radius 4
    \item \textbf{Quadratic in $z$}: Can use quadratic formula
    \item \textbf{Completing the square}: $1+z+z^2 = (z+\frac{1}{2})^2 + \frac{3}{4}$
    \item \textbf{Optimization}: Need to maximize one component (imaginary part)
    \item \textbf{Geometric interpretation}: Locus forms an ellipse in complex plane
\end{itemize}

\subsection*{H. Strategic Orientation}
\begin{itemize}[leftmargin=*]
    \item \textbf{Primary Strategy}: Use quadratic formula and optimization
    \item \textbf{Key Insight}: Maximum occurs when $1+z+z^2 = -4$ (most negative real)
    \item \textbf{Expected Difficulty}: Challenging - requires complex number manipulation and optimization reasoning
    \item \textbf{Time Pressure}: Should be completed in 4-5 minutes
\end{itemize}
\end{problembox}

\begin{strategybox}
\section*{2. Strategic Framework Application}

\subsection*{A. Psychological Strategies}
\begin{itemize}[leftmargin=*]
    \item \textbf{Confidence Calibration}: 
    \begin{itemize}
        \item Upper-mid difficulty problem
        \item Requires solid understanding of complex numbers
        \item Multiple approaches available (algebraic, geometric)
    \end{itemize}
    \item \textbf{Pattern Recognition}:
    \begin{itemize}
        \item Recognize quadratic form in $z$
        \item Completing the square technique
        \item Optimization of one variable subject to constraint
    \end{itemize}
    \item \textbf{Problem Familiarity}:
    \begin{itemize}
        \item Similar to constrained optimization problems
        \item Related to locus problems in complex plane
        \item Connection to AM-GM type inequalities
    \end{itemize}
\end{itemize}

\subsection*{B. Investigation Strategies}
\begin{itemize}[leftmargin=*]
    \item \textbf{Initial Exploration}:
    \begin{itemize}
        \item Let $z = a + bi$ and expand $|1+z+z^2|$
        \item Apply quadratic formula: $z = \frac{-1 \pm \sqrt{4c-3}}{2}$ where $|c| = 4$
        \item Consider when imaginary part is maximized
    \end{itemize}
    \item \textbf{Key Observations}:
    \begin{itemize}
        \item $1+z+z^2 = c$ where $|c| = 4$
        \item $c$ lies on circle of radius 4 centered at origin
        \item Imaginary part of $z$ depends on square root of $c$
    \end{itemize}
    \item \textbf{Optimization Reasoning}:
    \begin{itemize}
        \item To maximize $\text{Im}(z)$, need to maximize $\text{Im}(\sqrt{4c-3})$
        \item When is $\sqrt{4c-3}$ most purely imaginary?
        \item Answer: When $4c-3$ is most negative (i.e., $c = -4$)
    \end{itemize}
\end{itemize}

\subsection*{C. Argument Construction}
\begin{itemize}[leftmargin=*]
    \item \textbf{Logical Structure (Quadratic Formula Method)}:
    \begin{enumerate}
        \item Apply quadratic formula to $z^2 + z + (1-c) = 0$
        \item Get $z = \frac{-1 \pm \sqrt{4c-3}}{2}$
        \item To maximize $\text{Im}(z)$, need to maximize $\text{Im}(\sqrt{4c-3})$
        \item This occurs when $c = -4$ (on the circle $|c| = 4$)
        \item Compute: $z = \frac{-1 \pm \sqrt{-19}}{2} = \frac{-1 \pm i\sqrt{19}}{2}$
        \item Maximum imaginary part: $\frac{\sqrt{19}}{2}$
        \item Therefore: $m = 19$, $n = 2$, and $m + n = 21$
    \end{enumerate}
    \item \textbf{Key Deductions}:
    \begin{itemize}
        \item $|c| = 4$ means $c$ is on circle of radius 4
        \item $4c - 3$ is also on a circle (centered at $-3$, radius 16)
        \item Most negative point is $4c - 3 = -19$
        \item $\sqrt{-19} = i\sqrt{19}$ is purely imaginary
    \end{itemize}
\end{itemize}

\subsection*{D. Cross-Domain Translation}
\begin{itemize}[leftmargin=*]
    \item \textbf{Algebraic}: Solve quadratic equation with parameter
    \item \textbf{Geometric}: Locus of $z$ forms ellipse, find highest point
    \item \textbf{Optimization}: Constrained optimization using calculus or geometric reasoning
    \item \textbf{Complex Analysis}: Magnitude and argument manipulation
\end{itemize}
\end{strategybox}

\begin{tacticalbox}
\section*{3. Tactical Methodology Selection}

\subsection*{A. Core Mathematical Tactics}
\begin{itemize}[leftmargin=*]
    \item \textbf{Quadratic Formula}: Solve for $z$ in terms of parameter $c$
    \item \textbf{Parametrization}: Let $1+z+z^2 = c$ with $|c| = 4$
    \item \textbf{Completing the Square}: Express as $(z+\frac{1}{2})^2 + \frac{3}{4}$
    \item \textbf{Optimization}: Maximize imaginary component
    \item \textbf{Geometric Reasoning}: Visualize locus in complex plane
\end{itemize}

\subsection*{B. Domain-Specific Tactics}
\begin{itemize}[leftmargin=*]
    \item \textbf{Complex Numbers}:
    \begin{itemize}
        \item $z = a + bi$ substitution
        \item Magnitude: $|w| = \sqrt{(\text{Re}(w))^2 + (\text{Im}(w))^2}$
        \item Square root: $\sqrt{re^{i\theta}} = \sqrt{r}e^{i\theta/2}$
    \end{itemize}
    \item \textbf{Optimization}:
    \begin{itemize}
        \item Identify which parameter to vary ($c$ on circle)
        \item Determine when objective function is maximized
        \item Geometric interpretation helps intuition
    \end{itemize}
    \item \textbf{Algebraic Manipulation}:
    \begin{itemize}
        \item Expand expressions carefully
        \item Factor quadratics when possible
        \item Simplify square roots
    \end{itemize}
\end{itemize}

\subsection*{C. Tactical Priorities}
\begin{enumerate}[leftmargin=*]
    \item Set up parametrization: $1+z+z^2 = c$ with $|c| = 4$ (60 seconds)
    \item Apply quadratic formula (60 seconds)
    \item Identify optimization condition: $c = -4$ (90 seconds)
    \item Calculate maximum imaginary part (60 seconds)
    \item Find $m + n$ (30 seconds)
\end{enumerate}
\end{tacticalbox}

\begin{conversionbox}
\section*{4. Conversion Techniques Application}

\subsection*{A. Method 1: Quadratic Formula (Most Direct)}
\begin{itemize}[leftmargin=*]
    \item Let $1+z+z^2 = c$ where $|c| = 4$
    \item Rearrange: $z^2 + z + (1-c) = 0$
    \item Apply quadratic formula:
    \[z = \frac{-1 \pm \sqrt{1 - 4(1-c)}}{2} = \frac{-1 \pm \sqrt{4c-3}}{2}\]
    \item The imaginary part of $z$ is:
    \[\text{Im}(z) = \frac{\text{Im}(\sqrt{4c-3})}{2}\]
    \item To maximize $\text{Im}(\sqrt{4c-3})$, we need $4c-3$ to have:
    \begin{itemize}
        \item Maximum magnitude (for largest $\sqrt{r}$)
        \item Argument that gives maximum $\sin(\theta/2)$
    \end{itemize}
    \item Since $|c| = 4$, the point $c$ lies on a circle of radius 4
    \item Therefore $4c - 3$ lies on a circle centered at $-3$ with radius $16$
    \item The leftmost point of this circle is $-3 - 16 = -19$
    \item So the optimal value is $4c - 3 = -19$, which gives $c = -4$
    \item Then:
    \[z = \frac{-1 \pm \sqrt{-19}}{2} = \frac{-1 \pm i\sqrt{19}}{2}\]
    \item Maximum imaginary part: $\frac{\sqrt{19}}{2}$
    \item Therefore: $m = 19$, $n = 2$, $\gcd(19,2) = 1$ $\checkmark$
    \item Answer: $m + n = 21$
\end{itemize}

\subsection*{B. Method 2: Substitution and Optimization}
\begin{itemize}[leftmargin=*]
    \item Let $z = a + bi$
    \item Expand:
    \begin{align*}
    1 + z + z^2 &= 1 + (a+bi) + (a+bi)^2\\
    &= 1 + a + a^2 - b^2 + (b + 2ab)i\\
    &= (1 + a + a^2 - b^2) + (b + 2ab)i
    \end{align*}
    \item Apply magnitude constraint:
    \[(1 + a + a^2 - b^2)^2 + (b + 2ab)^2 = 16\]
    \item Expand the second term:
    \[(b + 2ab)^2 = b^2(1 + 2a)^2 = b^2(1 + 4a + 4a^2)\]
    \item Let $p = b^2$ and $q = 1 + a + a^2$:
    \[(q - p)^2 + p(1 + 4a + 4a^2) = 16\]
    \item Note: $1 + 4a + 4a^2 = 4(a^2 + a + \frac{1}{4}) = 4q - 3$
    \item So: $(q - p)^2 + p(4q - 3) = 16$
    \item Expand: $p^2 - 2pq + q^2 + 4pq - 3p = 16$
    \item Simplify: $p^2 + (2q - 3)p + (q^2 - 16) = 0$
    \item Solve for $p$:
    \[p = \frac{-(2q-3) \pm \sqrt{(2q-3)^2 - 4(q^2-16)}}{2}\]
    \[= \frac{3-2q \pm \sqrt{4q^2 - 12q + 9 - 4q^2 + 64}}{2}\]
    \[= \frac{3-2q \pm \sqrt{-12q + 73}}{2}\]
    \item To maximize $b = \sqrt{p}$, maximize $p$
    \item Take the $+$ sign and minimize $q$
    \item Note: $q = 1 + a + a^2 = (a + \frac{1}{2})^2 + \frac{3}{4} \geq \frac{3}{4}$
    \item Minimum $q = \frac{3}{4}$ when $a = -\frac{1}{2}$
    \item Then:
    \[p = \frac{3 - 2 \cdot \frac{3}{4} + \sqrt{-12 \cdot \frac{3}{4} + 73}}{2} = \frac{3 - \frac{3}{2} + \sqrt{64}}{2} = \frac{\frac{3}{2} + 8}{2} = \frac{19}{4}\]
    \item Therefore: $b = \sqrt{\frac{19}{4}} = \frac{\sqrt{19}}{2}$
    \item Answer: $m + n = 19 + 2 = 21$
\end{itemize}

\subsection*{C. Method 3: Geometric/Completing the Square}
\begin{itemize}[leftmargin=*]
    \item Complete the square:
    \[1 + z + z^2 = (z + \frac{1}{2})^2 + \frac{3}{4}\]
    \item So:
    \[\left|(z + \frac{1}{2})^2 + \frac{3}{4}\right| = 4\]
    \item Let $u = (z + \frac{1}{2})^2$, then $|u + \frac{3}{4}| = 4$
    \item This means $u$ lies on a circle centered at $-\frac{3}{4}$ with radius 4
    \item Now let $v = z + \frac{1}{2}$, so $v^2 = u$
    \item Since adding $\frac{1}{2}$ doesn't change imaginary part, we want to maximize $\text{Im}(v)$
    \item The point on the $u$-circle with largest magnitude is $-\frac{3}{4} - 4 = -\frac{19}{4}$
    \item So $v^2 = -\frac{19}{4}$, which gives $v = \pm \frac{i\sqrt{19}}{2}$
    \item Taking positive imaginary part: $v = \frac{i\sqrt{19}}{2}$
    \item Therefore: $z = v - \frac{1}{2} = -\frac{1}{2} + \frac{i\sqrt{19}}{2}$
    \item Maximum imaginary part: $\frac{\sqrt{19}}{2}$
    \item Answer: $m + n = 19 + 2 = 21$
\end{itemize}

\subsection*{D. Method 4: Factoring Approach}
\begin{itemize}[leftmargin=*]
    \item Factor: $1 + z + z^2 = (z - \omega)(z - \overline{\omega})$ where $\omega = e^{2\pi i/3} = \frac{-1 + i\sqrt{3}}{2}$
    \item So:
    \[\left|z - \omega\right| \cdot \left|z - \overline{\omega}\right| = 4\]
    \item Let $z = a + bi$. To maximize $b$, set $a = -\frac{1}{2}$ (the real part of $\omega$)
    \item Then:
    \[\left|bi + \frac{i\sqrt{3}}{2}\right| \cdot \left|bi - \frac{i\sqrt{3}}{2}\right| = 4\]
    \[\left|b + \frac{\sqrt{3}}{2}\right| \cdot \left|b - \frac{\sqrt{3}}{2}\right| = 4\]
    \[\left(b + \frac{\sqrt{3}}{2}\right)\left(b - \frac{\sqrt{3}}{2}\right) = 4\]
    \[b^2 - \frac{3}{4} = 4\]
    \[b^2 = \frac{19}{4}\]
    \[b = \frac{\sqrt{19}}{2}\]
    \item Answer: $m + n = 19 + 2 = 21$
\end{itemize}
\end{conversionbox}

\begin{verificationbox}
\section*{5. Verification Strategy Design}

\subsection*{A. Verification Level}
\begin{itemize}[leftmargin=*]
    \item \textbf{Selected Level}: Level 4 - Direct Substitution
    \item \textbf{Reasoning}: Can verify by checking constraint
    \item \textbf{Time Budget}: 60 seconds
\end{itemize}

\subsection*{B. Verification Checklist}
\begin{itemize}[leftmargin=*]
    \item \textbf{Check solution}: $z = -\frac{1}{2} + \frac{i\sqrt{19}}{2}$
    \begin{itemize}
        \item Compute $z^2 = \frac{1}{4} - \frac{19}{4} - \frac{i\sqrt{19}}{2} = -\frac{18}{4} - \frac{i\sqrt{19}}{2} = -\frac{9}{2} - \frac{i\sqrt{19}}{2}$
        \item Compute $1 + z + z^2 = 1 - \frac{1}{2} + \frac{i\sqrt{19}}{2} - \frac{9}{2} - \frac{i\sqrt{19}}{2}$
        \item $= 1 - 5 = -4$ $\checkmark$
        \item Magnitude: $|-4| = 4$ $\checkmark$
    \end{itemize}
    \item \textbf{Check form}:
    \begin{itemize}
        \item Imaginary part: $\frac{\sqrt{19}}{2}$ $\checkmark$
        \item $\gcd(19, 2) = 1$ $\checkmark$
    \end{itemize}
    \item \textbf{Check answer}:
    \begin{itemize}
        \item $m + n = 19 + 2 = 21$ $\checkmark$
        \item Matches choice (B) $\checkmark$
    \end{itemize}
\end{itemize}

\subsection*{C. Alternative Verification Methods}
\begin{itemize}[leftmargin=*]
    \item \textbf{Method 1}: Verify that $b = \frac{\sqrt{19}}{2}$ is indeed maximum (check derivative or boundary)
    \item \textbf{Method 2}: Check that other values of $c$ give smaller imaginary parts
    \item \textbf{Method 3}: Verify geometric interpretation (ellipse, highest point)
\end{itemize}

\subsection*{D. Answer Choice Validation}
\begin{itemize}[leftmargin=*]
    \item Our answer: 21
    \item Matches choice: \textbf{(B)}
    \item Why not the others?
    \begin{itemize}
        \item Answer must be computed, not guessed
        \item Other values don't satisfy the constraint
    \end{itemize}
\end{itemize}
\end{verificationbox}

\begin{roadmapbox}
\section*{6. Solution Roadmap}

\subsection*{A. High-Level Solution Path}
\begin{enumerate}[leftmargin=*]
    \item \textbf{Parametrize}: Let $1+z+z^2 = c$ with $|c| = 4$
    \item \textbf{Apply quadratic formula}: $z = \frac{-1 \pm \sqrt{4c-3}}{2}$
    \item \textbf{Optimize}: Maximize $\text{Im}(\sqrt{4c-3})$ by choosing $c = -4$
    \item \textbf{Calculate}: $z = \frac{-1 \pm i\sqrt{19}}{2}$
    \item \textbf{Extract answer}: $\text{Im}(z)_{\max} = \frac{\sqrt{19}}{2}$, so $m+n = 21$
\end{enumerate}

\subsection*{B. Detailed Solution Steps}

\textbf{Step 1}: Set up the problem.
\begin{itemize}
    \item Given: $|1+z+z^2| = 4$
    \item Let $1+z+z^2 = c$ where $c$ is some complex number with $|c| = 4$
\end{itemize}

\textbf{Step 2}: Solve for $z$ using quadratic formula.
\begin{itemize}
    \item Rearrange: $z^2 + z + (1-c) = 0$
    \item Apply quadratic formula:
    \[z = \frac{-1 \pm \sqrt{1-4(1-c)}}{2} = \frac{-1 \pm \sqrt{4c-3}}{2}\]
\end{itemize}

\textbf{Step 3}: Identify optimization condition.
\begin{itemize}
    \item Imaginary part of $z$ is $\frac{\text{Im}(\sqrt{4c-3})}{2}$
    \item To maximize this, need to maximize $\text{Im}(\sqrt{4c-3})$
    \item Since $|c| = 4$, $c$ lies on circle of radius 4
    \item Then $4c-3$ lies on circle centered at $-3$ with radius 16
    \item For $\sqrt{w}$ with $w = re^{i\theta}$, we have $\sqrt{w} = \sqrt{r}e^{i\theta/2}$
    \item So $\text{Im}(\sqrt{w}) = \sqrt{r}\sin(\theta/2)$
    \item This is maximized when both $r$ and $\sin(\theta/2)$ are large
    \item The point on the circle with largest $r$ and $\theta = \pi$ is $w = -19$
    \item Therefore $4c - 3 = -19$, giving $c = -4$
\end{itemize}

\textbf{Step 4}: Calculate maximum imaginary part.
\begin{itemize}
    \item With $c = -4$:
    \[z = \frac{-1 \pm \sqrt{4(-4)-3}}{2} = \frac{-1 \pm \sqrt{-19}}{2} = \frac{-1 \pm i\sqrt{19}}{2}\]
    \item Taking positive imaginary part: $\text{Im}(z) = \frac{\sqrt{19}}{2}$
\end{itemize}

\textbf{Step 5}: Find $m + n$.
\begin{itemize}
    \item $\frac{\sqrt{19}}{2}$ is in the form $\frac{\sqrt{m}}{n}$ with $m = 19$, $n = 2$
    \item Check: $\gcd(19, 2) = 1$ $\checkmark$
    \item Answer: $m + n = 19 + 2 = 21$
\end{itemize}

\subsection*{C. Key Decision Points}
\begin{itemize}[leftmargin=*]
    \item \textbf{Decision 1}: Parametrize with $c$ on circle
    \item \textbf{Decision 2}: Use quadratic formula
    \item \textbf{Decision 3}: Recognize that $c = -4$ maximizes imaginary part
    \item \textbf{Decision 4}: Correctly compute square root of negative number
\end{itemize}

\subsection*{D. Time Management}
\begin{itemize}[leftmargin=*]
    \item Understanding problem: 30 seconds
    \item Setting up parametrization: 60 seconds
    \item Applying quadratic formula: 60 seconds
    \item Optimization reasoning: 90 seconds
    \item Calculation: 60 seconds
    \item Verification: 60 seconds
    \item \textbf{Total}: 5 minutes
\end{itemize}
\end{roadmapbox}

\begin{competitionbox}
\section*{7. Competition Strategy Integration}

\subsection*{A. Problem Triage}
\begin{itemize}[leftmargin=*]
    \item \textbf{Difficulty Assessment}: Upper-medium (Problem 16/25)
    \item \textbf{Confidence Level}: Moderate to high if comfortable with complex numbers
    \item \textbf{Decision}: Skip if unfamiliar with complex numbers; attempt if comfortable
\end{itemize}

\subsection*{B. Time Allocation}
\begin{itemize}[leftmargin=*]
    \item \textbf{Target Time}: 4-5 minutes
    \item \textbf{Maximum Time}: 6 minutes
    \item \textbf{Abandonment Criteria}: If no progress after 4 minutes, move on
\end{itemize}

\subsection*{C. Solution Format}
\begin{itemize}[leftmargin=*]
    \item Multiple choice: Select \textbf{(B) 21}
    \item Verification: Check constraint with computed $z$
    \item Bubble answer sheet carefully
\end{itemize}

\subsection*{D. Strategic Considerations}
\begin{itemize}[leftmargin=*]
    \item \textbf{Pattern Recognition}: If you recognize the quadratic formula approach, this is faster
    \item \textbf{Elimination}: Cannot eliminate choices without solving
    \item \textbf{Key insight}: $c = -4$ gives maximum - this is the crucial step
    \item \textbf{Verification}: Always check $\gcd(m,n) = 1$
\end{itemize}
\end{competitionbox}

\begin{pitfallbox}
\section*{Common Pitfalls and Warnings}

\subsection*{A. Conceptual Errors}
\begin{itemize}[leftmargin=*]
    \item \textbf{Wrong parametrization}: Not recognizing that $c$ lies on circle
    \item \textbf{Optimization error}: Thinking maximum occurs at $c = 4$ instead of $c = -4$
    \item \textbf{Square root confusion}: $\sqrt{-19} = i\sqrt{19}$, not $-\sqrt{19}$
    \item \textbf{Missing imaginary unit}: Forgetting that square root of negative is imaginary
\end{itemize}

\subsection*{B. Computational Errors}
\begin{itemize}[leftmargin=*]
    \item \textbf{Quadratic formula}: Sign errors or arithmetic mistakes
    \item \textbf{Completing square}: $(z + \frac{1}{2})^2 = z^2 + z + \frac{1}{4}$, not $z^2 + z + 1$
    \item \textbf{GCD check}: Must verify $\gcd(19, 2) = 1$ before adding
    \item \textbf{Arithmetic}: $19 + 2 = 21$, not 20 or 22
\end{itemize}

\subsection*{C. Strategic Errors}
\begin{itemize}[leftmargin=*]
    \item \textbf{Overcomplicated algebra}: Using substitution method is much longer
    \item \textbf{Missing geometric intuition}: Not seeing why $c = -4$ is optimal
    \item \textbf{Time management}: Spending too long on algebraic manipulation
\end{itemize}

\subsection*{D. Warning Signs}
\begin{itemize}[leftmargin=*]
    \item If you get imaginary part with wrong denominator, recheck calculation
    \item If answer isn't among choices, revisit optimization step
    \item If $m$ and $n$ aren't coprime, you made an error
\end{itemize}
\end{pitfallbox}

\begin{solutionbox}
\subsection*{A. Main Solution (Quadratic Formula)}
\begin{enumerate}[leftmargin=*]
    \item \textbf{Parametrize the constraint}:
    \begin{itemize}
        \item Let $1 + z + z^2 = c$ where $|c| = 4$
        \item This means $c$ lies on a circle of radius 4 in the complex plane
    \end{itemize}
    
    \item \textbf{Solve for $z$}:
    \begin{itemize}
        \item Rearrange: $z^2 + z + (1-c) = 0$
        \item Quadratic formula:
        \[z = \frac{-1 \pm \sqrt{1 - 4(1-c)}}{2} = \frac{-1 \pm \sqrt{4c-3}}{2}\]
    \end{itemize}
    
    \item \textbf{Optimize for maximum imaginary part}:
    \begin{itemize}
        \item $\text{Im}(z) = \frac{\text{Im}(\sqrt{4c-3})}{2}$
        \item Since $|c| = 4$, we have $c = 4e^{i\alpha}$ for some angle $\alpha$
        \item Then $4c - 3 = 16e^{i\alpha} - 3$
        \item This traces a circle centered at $-3$ with radius 16
        \item The leftmost point is $-3 - 16 = -19$ (when $\alpha = \pi$, i.e., $c = -4$)
        \item For $w = re^{i\theta}$, we have $\sqrt{w} = \sqrt{r}e^{i\theta/2}$
        \item With $w = -19 = 19e^{i\pi}$: $\sqrt{-19} = \sqrt{19}e^{i\pi/2} = i\sqrt{19}$
    \end{itemize}
    
    \item \textbf{Calculate maximum}:
    \begin{align*}
    z &= \frac{-1 \pm i\sqrt{19}}{2}\\
    \text{Im}(z)_{\max} &= \frac{\sqrt{19}}{2}
    \end{align*}
    
    \item \textbf{Find answer}:
    \begin{itemize}
        \item $m = 19$, $n = 2$
        \item $\gcd(19, 2) = 1$ $\checkmark$
        \item $m + n = 21$
    \end{itemize}
\end{enumerate}

\subsection*{B. Verification}
\begin{itemize}[leftmargin=*]
    \item Check: $z = -\frac{1}{2} + \frac{i\sqrt{19}}{2}$
    \item Compute:
    \begin{align*}
    z^2 &= \frac{1}{4} - \frac{19}{4} - \frac{i\sqrt{19}}{2} = -\frac{9}{2} - \frac{i\sqrt{19}}{2}\\
    1 + z + z^2 &= 1 - \frac{1}{2} + \frac{i\sqrt{19}}{2} - \frac{9}{2} - \frac{i\sqrt{19}}{2}\\
    &= 1 - 5 = -4
    \end{align*}
    \item $|-4| = 4$ $\checkmark$
\end{itemize}
\end{solutionbox}

\begin{keyinsightbox}
\textbf{Key Insight}: The crucial observation is that to maximize the imaginary part of $z = \frac{-1 \pm \sqrt{4c-3}}{2}$, we need to maximize $\text{Im}(\sqrt{4c-3})$.

For a complex number $w = re^{i\theta}$:
\[\text{Im}(\sqrt{w}) = \text{Im}(\sqrt{r}e^{i\theta/2}) = \sqrt{r}\sin(\theta/2)\]

This is maximized when:
\begin{itemize}
    \item $r$ is as large as possible
    \item $\sin(\theta/2)$ is as large as possible
\end{itemize}

Since $4c - 3$ lies on a circle centered at $-3$ with radius 16, the point with:
\begin{itemize}
    \item Largest magnitude: $|-3 - 16| = 19$
    \item Argument $\theta = \pi$ (negative real axis): $\sin(\pi/2) = 1$
\end{itemize}

is $w = -19$, which occurs when $c = -4$.

\textbf{General Pattern}: For optimization problems involving $\sqrt{f(c)}$ where $c$ is on a circle:
\begin{itemize}
    \item Consider where $f(c)$ is located
    \item Find the point that maximizes both magnitude and favorable argument
    \item For imaginary part, want argument $\pi$ (negative real)
    \item For real part, want argument $0$ (positive real)
\end{itemize}

\textbf{Why $c = -4$ works}: 
\begin{itemize}
    \item $4c - 3 = -19$ is the most negative real number possible
    \item $\sqrt{-19} = i\sqrt{19}$ is purely imaginary with maximum magnitude
    \item This gives the largest positive imaginary part for $z$
\end{itemize}
\end{keyinsightbox}

\section*{Final Answer}

The maximum imaginary part of $z$ is $\frac{\sqrt{19}}{2}$, so:

\[\boxed{\textbf{(B) } 21}\]

\textbf{Solution Summary}:
\begin{itemize}
    \item Parametrize: $1+z+z^2 = c$ with $|c| = 4$
    \item Solve: $z = \frac{-1 \pm \sqrt{4c-3}}{2}$
    \item Optimize: Maximum imaginary part when $c = -4$
    \item Calculate: $z = \frac{-1 \pm i\sqrt{19}}{2}$, so $\text{Im}(z)_{\max} = \frac{\sqrt{19}}{2}$
    \item Answer: $m + n = 19 + 2 = 21$
\end{itemize}

\section*{Confidence Assessment}
\begin{itemize}[leftmargin=*]
    \item \textbf{Confidence Level}: 95\%
    \begin{itemize}
        \item Standard quadratic formula approach
        \item Clear optimization condition ($c = -4$)
        \item Verified by direct substitution
        \item Answer matches expected form
    \end{itemize}
    \item \textbf{Verification Applied}: Level 4 - Direct Substitution
    \begin{itemize}
        \item Checked that $z = -\frac{1}{2} + \frac{i\sqrt{19}}{2}$ satisfies constraint
        \item Confirmed $\gcd(19, 2) = 1$
        \item Verified arithmetic: $19 + 2 = 21$
    \end{itemize}
    \item \textbf{Time-Confidence Trade-off}: High confidence in 4-5 minutes
\end{itemize}

\end{document}

